\section{Parameter estimates after the local projection}
\label{sec:projected-parameter-prices}
We transfer the estimates of Section~\ref{sec:sector-analytic-prices}
to the actual ordinary and forced arrays after localization. All
finite tests below use the common parameter family, and every local
cutoff is applied to its original fixed support. We retain the
compatibility and local-generation identities specified in
Section~\ref{sec:five-period-localization}.

The ordinary budget below is anchored: at site $o$ it sums the
nonnegative derivative prices of entries whose assigned anchor is $o$,
with their exponential support weights. The forced budget sums all
supports containing the site. The sector bounds imply the anchored
ordinary bounds by restricting their support sums to the assigned anchors.

Let $P$ be the outgoing derivative budget, with derivative order
$j$, spatial weight power $m$, radius $B$, threshold $\theta$ and
support exponent $b$. Put $b_*=b+3$. Let
$y=\beta_f(x)$ and $R=y^{1/256}$, as before. Assume the path
rows have their purchased strip and coefficient bounds, their
regular local readings are smooth on healthy supports, and the
support reserve fits inside half the path strip. Their ordinary
rooted allowance, before the finite orientation average, is at
most $y^{-7/16}$. These are the path inputs to the present
argument. The parameter-dependent action and quadratic rows
are multiplied by $\ell(t)^2$, while the other two row types
are unchanged. Since $|\ell^2|_R\le(1+\delta)^2\le4$, their
finite-test allowance at radius $R/2$ is at most four times
the original allowance.

Let $T_P>0$ be the fixed orientation and derivative allocation
cost. Combining the path terms and the sector logarithmic terms,
and applying that allocation, gives the price
\begin{equation}\label{eq:projected-initial-price}
                    W(x)=T_P\{y^{-3/8}+4y^{-7/16}\}.
\end{equation}
For every finite test $\lambda$, the ordinary family has rooted
derivative price $M_{R/2}(\lambda)W(x)$ at exponent $b_*$.
A fixed finite short-support panel may be removed from this
family with the same bound: keep the original nonnegative price
on the retained addresses and put zero on the deleted ones.
Neither the parameter dictionary nor its support sets change.
This observation includes every finite truncation used in the
short-density construction, with no dependence of $W$ on the
truncation order.

\subsection{Allocating and testing the auxiliary densities}
For the five auxiliary periods $\Lambda_i$ of
Section~\ref{sec:five-period-localization}, form the same
finite-panel local densities $f_i(q,U)$. The nonnegative
allocation coefficients are bounded by one, sum to one over
$q$, and only assign an entry to anchors in its prescribed
neighbourhood. Write $s$ for the radius of that neighbourhood
and $b_s=(2s+1)^4$.

For a tested density at a fixed site $q$, take absolute values
in its allocation formula. Each nonzero entry is assigned to
an anchor within $N_s(q)$. Summing first over those anchors
and then using the rooted bound gives
\begin{equation}\label{eq:projected-local-test}
  \left|\sum_h a_h f_i(t_h,q,U)\right|
                         \le M_{R/2}(\lambda)b_sW(x)
\end{equation}
on globally healthy $U$. The same estimate holds for the
masked panel on an arbitrary field: outside its healthy support
an entry is zero. Here the mask is independent of the parameter,
so testing commutes with masking.

The total auxiliary density is exactly the sum of the panel
entries, since the allocations sum to one. Summing the derivative
prices over anchors therefore bounds its derivatives by
\[ |\Lambda_i|W(x)M_{R/2}(\lambda). \]
Subtracting its value at the
identity field costs at most twice that amount, uniformly through
derivative order $j$; positive derivatives of the subtracted
constant vanish. Thus the two probe prices needed below are
\begin{equation}\label{eq:projected-probe-prices}
                     p_i=2|\Lambda_i|W(x),\qquad q_i=b_sW(x).
\end{equation}
The first bounds derivatives of the centered total density;
the second bounds each uncentered local density on healthy fields.
The identity field is healthy because its plaquette holonomies
are all the identity and $\theta>0$.

\subsection{The ordinary output price}
Let $g(t,o,U)$ be the five-period local density and $d(t,o,U)$
its discrepancy. Put
\[
 \mathcal O_z(p)=\sum_{i=0}^4|c_i|\,2^j(1+z)^{mj}p_i,
 \qquad C_z\ge\sup_o c_B(N_z(o)).
\]
All operations in the formula for $g$ are finite linear
operations on the densities. Hence a finite test commutes with
that formula. Apply \eqref{eq:five-period-derivative-cost} to
the tested auxiliary total densities. The result is the regular
support price $M_{R/2}(\lambda)\mathcal O_z(p)$ at each anchor.
Multiplying by the support weight gives
$M_{R/2}(\lambda)\mathcal O_z(p)e^{b_*C_z}$.
The existing complementary panel has price
$M_{R/2}(\lambda)W(x)$. After the orientation allocation,
the total ordinary output price is therefore
\begin{equation}\label{eq:projected-ordinary-price}
       A(x)=W(x)+T_P\mathcal O_z(p)e^{b_*C_z}.
\end{equation}
This is a bound on the full derivative budget on each healthy
support, not just on the sum of the array values.

In the selected geometry, $s,z$ and all five auxiliary volumes
have upper exponent zero; the cover bound $C_z$ is fixed by
the halo and support construction. The derivative order and
spatial weight power are fixed before the scale is increased.
Consequently $A$ has upper exponent $-3/8$, and
\begin{equation}\label{eq:projected-broad-bound}
                                A(x)\le y^{-1/4}
\end{equation}
eventually. The strict exponent gap absorbs every fixed
orientation constant and every factor of upper exponent zero.
This is the broad ordinary bound used in the finite-test
sharpening of Section~\ref{sec:background-test-prices}; its
radius is $R/2=y^{1/256}/2$.

\subsection{A common series majorant for the discrepancy}
A discrepancy value is the difference between the allocated
masked panel, its identity-field constant, and the shifted
masked auxiliary sum. Applying \eqref{eq:projected-local-test}
to the first two terms gives $2b_sW(x)M_{R/2}(\lambda)$.
For the last term, each centered auxiliary local value is bounded
by $2q_iM_{R/2}(\lambda)$ when its mask is one, since the folded
field is then globally healthy. It is zero when its mask is zero.
Hence, for every field $U$ and anchor $o$,
\[
 \left|\sum_h a_h d(t_h,o,U)\right|
 \le M_{R/2}(\lambda)
      \left\{2b_sW(x)+\sum_{i=0}^4|c_i||\Lambda_i|\,2q_i\right\}.
\]
Write $D$ for the expression in braces. The finite-test functional
argument gives coefficients $d_n(o,U)$ satisfying
$|d_n(o,U)|\le D(2/R)^n$ for every $U$, with the same $D$.
Therefore $D(2/R)^n$ is a common coefficientwise envelope; its
weighted sum at radius $R/4$ is $\sum_nD2^{-n}=2D$. Thus
\begin{equation}\label{eq:projected-discrepancy-value}
 |d(\cdot,o,U)|_{R/4}\le a(x),\qquad
 a(x)=2\left\{2b_sW(x)+\sum_{i=0}^4|c_i||\Lambda_i|\,2q_i\right\}.
\end{equation}
This bound is uniform in $U$, and $R\ge4$ eventually. In
particular there is one coefficient majorant for the discrepancy
at each fixed support. The sum over orientation copies uses the
normalized finite orientation average; for values it has norm
one. Assigning each copy its actual averaging weight and counting
the anchors whose radius-$2z$ support contains a root gives the
rooted price
\begin{equation}\label{eq:projected-discrepancy-root}
                H(x)=(4z+1)^4 a(x)e^{(b+2)C_{2z}},
                  \qquad C_{2z}\ge\sup_o c_B(N_{2z}(o)).
\end{equation}
This has upper exponent $-3/8$. The signs of the discrepancy
are retained in the functions; they disappear only in this
nonnegative bound.

\subsection{Fusing the discrepancy with the forced array}
Let $\gamma(x)$ be the forced series allowance supplied by
\eqref{eq:sector-forced-final}, and let $\phi(x)$ be the
prescribed outgoing forced allowance. Use the sector allowance
$\gamma$ at support exponent $b+2$, weakening its available
$b+3$ bound. The discrepancy allowance $H$ is also at exponent
$b+2$. The final forced bound will have exponent $b+1$, since
the connected-set estimate consumes the reserve of one. Both
allowances are evaluated at the
same scale and the same reference band. The numerical payment
chooses $\kappa\ge0$ so that, with
\[
                          I=e^\kappa\gamma+2H,
\]
the support-tree criterion at reserve one holds and
\begin{equation}\label{eq:projected-fusion-payment}
               H\le1,\qquad 4b_BI\le\tfrac12,
                              \qquad 2e^{-\kappa}I\le\phi.
\end{equation}
These are the inflation, degree and return inequalities of the
exponential payment argument in
Section~\ref{sec:nonlinear-payment-selection}. Its hypotheses
hold here: $H$ has negative upper exponent, while $\gamma$
retains the compulsory exponentially small mark. The prescribed
outgoing gain leaves the stated positive margin in that mark.
The selected payment is uniform over the coupling band.

Here is why those inequalities pay the analytic coefficients,
not only the real values. Give an old forced entry with series
majorant $k_l$ the inflated price $e^\kappa k_l$, and give a
discrepancy entry with series price $h_l$ the price $2h_l$.
A nonempty support and the rooted bound imply $h_l\le H\le1$.
The coefficient algebra therefore gives
\[
                 |e^{-d_l}-1|_{R/4}
                     \le e^{h_l}-1\le2h_l.
\]
The last inequality holds on $[0,1]$, for example by convexity
of the exponential between its endpoint values and $e-1<2$.
Every final forced component contains at least one old forced
entry. Its exact product of the old forced factors and the
exponential-minus-one discrepancy factors consequently has
series norm at most
\[
                 e^{-\kappa}\prod_{l\in D} F_l,
 \qquad F_l=\begin{cases}e^\kappa k_l,&l\text{ old forced},\\
                              2h_l,&l\text{ discrepancy}.
             \end{cases}
\]
Indeed each old forced factor contributes $e^{-\kappa}$;
retaining just one of them only increases the bound since
$\kappa\ge0$. The eligibility predicate for a component is
independent of $t$, so it preserves the same majorant. The
connected-set estimate, with rooted allowance $I$ and degree
ratio at most $1/2$, gives the final rooted allowance
$2e^{-\kappa}I\le\phi$. This proves a common coefficient
majorant at radius $R/4$ for all configurations and all entries
of the actual final forced array.

The exact additive decomposition of
Section~\ref{sec:five-period-localization} identifies these
ordinary and forced outputs with the selected one-step source.
Their dictionary is fixed by the common path and incoming
supports. The auxiliary test representations may be used outside
$[-1,1]$ to provide the real-node bounds, but on $[-1,1]$ they
agree with the given parameter family. Thus all the finite tests
and series recoveries above concern that same family. The ordinary
bound \eqref{eq:projected-broad-bound} and the forced common
majorant are precisely the two broad analytic inputs retained in
Section~\ref{sec:background-test-prices}. Construction of the
funded path and its real-sector hypotheses remains separate.
